\chapter{直角三角形与余弦}\label{ch:g8:cosine}

直角三角形与圆之间有两个美妙的事实：
直角三角形恰好内接于半圆，其斜边即为直径。
还有一个数，即角的\emph{余弦}，
刻画了每一个具有该角的直角三角形的形状——
它是第一个三角比，先于其兄弟正弦与正切
（见\cref{ch:g9:trig}）。

\section{直角三角形与它的圆}

\begin{theorem}[圆定理]\label{thm:g8:cosine:circle}
\begin{enumerate}
  \item 若三角形 $ABC$ 在 $A$ 处为直角，则 $A$ 位于
        以斜边 $[BC]$ 为直径的圆上。
  \item 反之，若 $A$ 位于直径为 $[BC]$ 的圆上（且
        $A \neq B, C$），则三角形 $ABC$ 在
        $A$ 处为直角。
\end{enumerate}
等价地：在直角三角形中，斜边的中点
到三个顶点的距离相等——由直角顶点引出的中线
长为斜边的一半。
\end{theorem}

\begin{proof}[1 的证明]
设 $O$ 为 $[BC]$ 的中点，$D$ 为 $A$ 关于
点 $O$ 的对称点。$ABDC$ 的两条对角线相交于其公共中点 $O$，
故 $ABDC$ 为平行四边形（见\cref{def:g7:central:parallelogram}）——
又 $A$ 处为直角，故为长方形。长方形的对角线
相等，故 $OA = \frac{AD}{2} = \frac{BC}{2}$，
即 $A$ 处于距 $\frac{BC}{2}$ 的位置，此距离量自 $O$，
亦即位于以 $[BC]$ 为直径的圆上。
（第 2 条只需将上述论证倒过来即可证明。）
\end{proof}

\begin{omfigure}
\begin{tikzpicture}[scale=1.15]
  \coordinate (B) at (-1.8,0);
  \coordinate (C) at (1.8,0);
  \coordinate (O) at (0,0);
  \coordinate (A) at (130:1.8);
  \coordinate (A2) at (55:1.8);
  \draw[thick] (O) circle (1.8);
  \draw[thick] (B) -- (C);
  \draw[thick, omDef] (B) -- (A) -- (C);
  \draw[thick, omExo] (B) -- (A2) -- (C);
  \draw[omThm, thick] (O) -- (A);
  \fill (B) circle (1.5pt) node[below left, font=\small] {$B$};
  \fill (C) circle (1.5pt) node[below right, font=\small] {$C$};
  \fill (O) circle (1.5pt) node[below, font=\small] {$O$};
  \fill (A) circle (1.5pt) node[above left, font=\small] {$A$};
  \fill (A2) circle (1.5pt) node[above right, font=\small] {$A'$};
  \draw[thin] (A) ++(0.14,-0.20) -- ++(0.14,0.10) -- ++(-0.14,0.20);
\end{tikzpicture}

{\small 无论 $A$ 位于圆上何处，角 $\widehat{BAC}$ 都
为直角——这是\cref{exo:g6:shapes:10}中的猜想，现已成为定理。
红色中线 $[OA]$ 是半径，即斜边的一半。}
\end{omfigure}

\begin{example}\label{ex:g8:cosine:circle}
一个三角形的斜边长为 $10$ cm。在不知道其他信息的情况下，
由直角顶点引出的中线长为 $5$ cm，且
该三角形的外接圆半径为 $5$ cm，圆心位于
斜边的中点。
\end{example}

\section{锐角的余弦}

\begin{definition}[余弦]\label{def:g8:cosine:def}
在直角三角形中，对锐角 $\theta$：
\[
\cos\theta = \frac{\text{邻边 } \theta}{\text{斜边}}
\]
——即与 $\theta$ 相邻的直角边除以斜边。该比值
只取决于角，而与三角形的大小无关：所有
具有相同锐角的直角三角形彼此互为放大图形，
而放大保持长度之比不变
（在\cref{ch:g8:midpoints}中已开其端；\cref{ch:g9:thales}
则完成之）。
\end{definition}

\begin{omfigure}
\begin{tikzpicture}[scale=1.05]
  \coordinate (B) at (0,0);
  \coordinate (A) at (4.6,0);
  \coordinate (C) at (4.6,2.3);
  \coordinate (A1) at (2.76,0);
  \coordinate (C1) at (2.76,1.38);
  \draw[thick] (B) -- (A) -- (C) -- cycle;
  \draw[omThm, thick] (A1) -- (C1);
  \draw[thin] (A) ++(-0.25,0) -- ++(0,0.25) -- ++(0.25,0);
  \draw[thin] (A1) ++(-0.22,0) -- ++(0,0.22) -- ++(0.22,0);
  \draw[omDef, ->] (0.85,0) arc (0:26.6:0.85);
  \node[omDef, font=\small] at (1.25,0.24) {$\theta$};
  \node[below left, font=\small] at (B) {$B$};
\end{tikzpicture}

{\small 两个直角三角形共用角 $\theta$：小三角形
是大三角形的缩小图形，因此
$\frac{\text{邻边}}{\text{斜边}}$ 对两者相同——
这一公共值即为 $\cos\theta$。}
\end{omfigure}

\begin{proposition}[初值与界]\label{prop:g8:cosine:bounds}
对每个锐角 $\theta$，有 $0 < \cos\theta < 1$，且余弦
\emph{随角增大而减小}：张角越大，
其余弦越小。特殊值：$\cos 0^\circ = 1$，$\cos 60^\circ = \frac12$，
$\cos 90^\circ = 0$（这些极端值分别对应
趋于扁平的三角形）。
\end{proposition}
\admitted

\begin{method}[利用余弦]\label{met:g8:cosine:use}
在直角三角形中，当已知量与所求量分别是
一个锐角、它的邻边以及斜边时：
\begin{enumerate}
  \item 写出定义：
        $\cos\theta = \frac{\text{邻边}}{\text{斜边}}$，并
        代入已知值；
  \item 解所得的小方程求出未知量（乘或除）；
  \item 求未知\emph{角}时，对算得的比值使用计算器的
        $\cos^{-1}$ 键；
  \item 检验结果是否合理：长度必须小于
        斜边；角必须严格介于 $0^\circ$ 与
        $90^\circ$ 之间。
\end{enumerate}
\end{method}

\begin{example}[求边长]\label{ex:g8:cosine:side}
一个 $6$ m 长的斜坡与水平地面成 $20^\circ$ 的角。
水平方向所走过的距离（与 $20^\circ$ 相邻的边，斜边
为 $6$ m）：
\[
\cos 20^\circ = \frac{d}{6}
\quad\Longrightarrow\quad
d = 6 \cos 20^\circ \approx 6 \times 0.940 \approx 5.6 \text{ 米}.
\]
\end{example}

\begin{example}[求角]\label{ex:g8:cosine:angle}
在直角三角形中，与角 $\theta$ 相邻的边长为
$3.5$，斜边为 $5$：
\[
\cos\theta = \frac{3.5}{5} = 0.7,
\qquad
\theta = \cos^{-1}(0.7) \approx 45.6^\circ .
\]
\end{example}

\section{练习}

\begin{exercise}[$\star$]\label{exo:g8:cosine:1}
一个直角三角形的斜边长为 $13$ cm。
求从直角顶点引出的中线长，
以及过三个顶点的圆的半径。
\end{exercise}

\begin{exercise}[$\star$]\label{exo:g8:cosine:2}
画一个直径为 $8$ cm 的圆，以 $[BC]$ 为直径，在圆上任取
一点 $A$，画出三角形 $ABC$。哪个角是直角？请引用
相应定理。量取 $AB$、$AC$ 的长，并验证勾股定理。
\end{exercise}

\begin{exercise}[$\star$]\label{exo:g8:cosine:3}
在直角三角形 $DEF$ 中，$D$ 为直角，写出与
角 $\widehat E$ 相邻的边、对边以及斜边。
用比表示 $\cos \widehat E$。
\end{exercise}

\begin{exercise}[$\star$]\label{exo:g8:cosine:4}
求下列未知量（$\cos 35^\circ \approx 0.819$，
$\cos 50^\circ \approx 0.643$）：
\begin{enumerate}
  \item 斜边为 $10$，角为 $35^\circ$：邻边是多少？
  \item 邻边为 $6$，角为 $50^\circ$：斜边是多少？
\end{enumerate}
\end{exercise}

\begin{exercise}[$\star$]\label{exo:g8:cosine:5}
在直角三角形中，与 $\theta$ 相邻的边长为 $8$，
斜边为 $10$。求 $\cos\theta$，再求 $\theta$
（$\cos^{-1}(0.8) \approx 36.9^\circ$）。
\end{exercise}

\begin{exercise}[$\star$]\label{exo:g8:cosine:6}
说明为什么对于直角三角形的锐角，
$\cos\theta$ 永远不可能等于 $1.2$。
\end{exercise}

\begin{exercise}[$\star\star$]\label{exo:g8:cosine:7}
一条 $25$ m 长的滑索从平台向地面下降，
与水平方向成 $12^\circ$ 的角（$\cos 12^\circ \approx
0.978$）。求它水平方向跨越的距离是多少？
\end{exercise}

\begin{exercise}[$\star\star$]\label{exo:g8:cosine:8}
不用计算器，将 $\cos 20^\circ$、$\cos 70^\circ$、
$\cos 45^\circ$ 按大小排序（利用\cref{prop:g8:cosine:bounds}）。
然后用计算器检验。
\end{exercise}

\begin{exercise}[$\star\star$]\label{exo:g8:cosine:9}
一架长为 $L$ 的梯子靠在墙上，梯子与
\emph{地面}之间的夹角为 $\theta$，梯脚距墙 $1.2$ m，
且 $\theta = 68^\circ$（$\cos 68^\circ \approx 0.375$）。求 $L$，
再求梯子所能到达的高度（利用勾股定理）。
\end{exercise}

\begin{exercise}[$\star\star$]\label{exo:g8:cosine:10}
$[BC]$ 是圆心为 $O$、半径为 $4.5$ cm 的圆的直径，
$A$ 是该圆上一点，且 $AB = 5.4$ cm。
\begin{enumerate}
  \item 为什么 $ABC$ 在 $A$ 处为直角？
  \item 求 $AC$，再求 $\cos \widehat B$。
\end{enumerate}
\end{exercise}

\begin{exercise}[$\star\star\star$]\label{exo:g8:cosine:11}
取边长为 $1$ 的等边三角形的一半（沿高剪开），
说明\cref{prop:g8:cosine:bounds}中的特殊值 $\cos 60^\circ = \frac12$。
（高的垂足落在底边的何处？
完整的三角函数版本见
\cref{ex:g9:trig:special}。）
\end{exercise}

\section{问题：欧几里得关系，以及
再证勾股定理}

\begin{problem}[{周末问题 --- 直角三角形的高：
$AB^2 = BH \times BC$，$AH^2 = BH \times HC$，勾股定理的第二种证明，
以及开平方根的机器}]
\label{pb:g8:cosine:1}
从直角三角形的直角处向斜边作垂线，
这条短线段 --- \emph{高} ---
满足一些非常有用的关系，以至于欧几里得把它们置于
其 \emph{《几何原本》} 的核心。本题中，余弦的定义性质
“比值只依赖于角”
（\cref{def:g8:cosine:def}）可证明所有这些关系；途中你将
用一条完全不同的路线重新证明勾股定理，
最后得到一台尺规作图机器，
可作出 $\sqrt2$、$\sqrt5$，以及 $\sqrt{n}$（对每个正整数 $n$）。

全文中，$ABC$ 是以 $A$ 为直角的三角形，$H$ 是
从 $A$ 向斜边 $[BC]$ 所作垂线的垂足，
从而 $H$ 位于 $B$ 与 $C$ 之间，且 $AH \perp BC$。

\medskip
\noindent\textbf{第一部分 --- 一个角，两个三角形。}
\begin{enumerate}
  \item 分别在 $ABH$ 与 $ABC$ 中应用三角形内角和
        （\cref{thm:g7:triangles:sum}），证明
        $\widehat{BAH} = \widehat{ACB}$：高把直角三角形分成两个较小的三角形，
        它们具有与原三角形
        \emph{相同的角}。
  \item 三角形 $ABC$（直角在 $A$ 处）与三角形 $ABH$
        （直角在 $H$ 处）有公共角 $\widehat B$。分别写出两者的
        $\cos\widehat B$，并由
        \cref{def:g8:cosine:def}，可得
        \[
        AB^2 = BH \times BC .
        \]
        （由相等的比得到相等的积，需将两边同乘以
        两个分母，
        依据\cref{thm:g8:equations:balance}）。
  \item 陈述并证明顶点 $C$ 处的对应关系：
        \[
        AC^2 = CH \times CB .
        \]
  \item 取 $3$--$4$--$5$ 三角形：$AB = 3$，$AC = 4$，
        $BC = 5$。由第 2、3 小题求 $BH$ 与 $CH$，
        并检验 $BH + HC = BC$。
  \item 一般地，将 $BH$ 与 $CH$ 表示为三边的分数，
        且仅涉及三边。垂足 $H$ 从何种意义上
        把斜边“按两直角边的
        平方成正比例”地分开？
\end{enumerate}

\medskip
\noindent\textbf{第二部分 --- 再证勾股定理与高。}
\begin{enumerate}[resume]
  \item 将第 2、3 小题中的关系式相加，并利用
        $BH + HC = BC$，得到
        \[
        AB^2 + AC^2 = BC^2 :
        \]
        这就是勾股定理（\cref{thm:g8:pythagoras:direct}）的
        再证明 --- 不借助任何面积拼图（可比较
        \cref{exo:g8:pythagoras:11}）。
  \item 在小三角形 $ABH$ 中应用勾股定理，写出
        $AH^2 = AB^2 - BH^2$，将 $AB^2$ 替换为
        $BH \times BC$，并提出公因式 $BH$，证明
        \emph{欧几里得高的关系}：
        \[
        AH^2 = BH \times HC .
        \]
  \item 用两种方式计算 $ABC$ 的面积 --- 先取两直角边分别作底与高，
        再取斜边作底
        （\cref{thm:g7:areas:triangle}） --- 并由此推出
        第三个关系：
        \[
        AB \times AC = BC \times AH .
        \]
  \item 回到 $3$--$4$--$5$ 三角形：利用第 8 小题求 $AH$，
        然后数值验证第 7 小题中的高的关系，
        所用的 $BH$ 与 $CH$ 的值取自
        第 4 小题。
  \item 某直角三角形的高的垂足把斜边分成
        两段，分别为 $BH = 2$ 厘米与 $HC = 8$ 厘米。
        求高 $AH$。
\end{enumerate}

\medskip
\noindent\textbf{第三部分 --- 开平方根的机器。}
作一条线段 $[BC]$，由首尾相接的两段组成：
$BH = p$，$HC = q$。以 $[BC]$ 为直径作半圆，
设 $A$ 为垂直于 $(BC)$ 且过 $H$ 的直线
与半圆的交点。
\begin{enumerate}[resume]
  \item 利用\cref{thm:g8:cosine:circle}，说明三角形
        $ABC$ 的直角在 $A$ 处 --- 因此可用第 7 小题，
        并得到
        \[
        AH^2 = p \times q .
        \]
        （线段 $AH$ 的长度称为\emph{几何
        平均数}\index{几何平均数}，它由 $p$ 与 $q$ 得到）。
  \item 取 $p = 1$，$q = 2$。$AH^2$ 等于多少？
        你刚用尺规作出的图形，构造出了
        \cref{ex:g8:pythagoras:diagonal}中的哪一个著名长度？
  \item 描述一种作图方法，可作出长为
        $\sqrt n$（$n \geq 1$ 为任意正整数）的线段，并指出
        选取哪组 $p$ 与 $q$ 可作出 $\sqrt 5$。
  \item 设 $O$ 为半圆的圆心。说明为什么
        $AH$ 不可能超过半径 $\frac{p + q}{2}$，
        并说明 $H$ 位于何处时二者相等。
  \item 推出不等式：对所有正数
        $p$ 与 $q$，
        \[
        p \times q \leq \left(\frac{p + q}{2}\right)^2,
        \]
        等号当且仅当 $p = q$ 时成立。然后完全不用几何，
        重新证明它：展开
        $\left(\frac{p+q}{2}\right)^2 - p q$，并识别
        一个著名的恒等式（见\cref{pb:g8:equations:1}）。
\end{enumerate}
\end{problem}
